%% sections/04_benchmarks.tex

\section{Numerical Benchmarks}
\label{sec:benchmark}

\subsection{Test Configuration}
\label{sec:testconfig}

We integrate three asteroids in purely two-body (Keplerian) dynamics
to isolate integrator performance from force-model effects.
Initial conditions in heliocentric ECLIPJ2000 are derived from
the osculating elements listed in Table~\ref{tab:elements}.
All integrations use $\mu = \mathrm{GMS}$ only.
Apophis is selected as a benchmark object for its 2029 Earth
close encounter, which makes it of broad observational relevance
\citep{Benson2022apophis, Wiegert2024, Broz2024apophis}.
CPU times are measured on Apple MacBook Air M-series, macOS,
clang C++23.

%% tables/tab_elements.tex

\begin{table*}
\centering
\caption{
  Osculating elements (heliocentric ECLIPJ2000, epoch J2000.0)
  of benchmark asteroids used in this work.
  \label{tab:elements}
}
\begin{tabular}{lccccccc}
\toprule
Body & $a$ (AU) & $e$ & $i$ (deg) & $\Omega$ (deg) & $\omega$ (deg) & $M$ (deg) \\
\midrule
(1) Ceres        & 2.7665 & 0.0784 & 10.583 &  80.494 &  73.923 &   6.177 \\
(99942) Apophis  & 0.9223 & 0.1914 &  3.331 & 204.657 & 126.065 & 232.181 \\
(3200) Phaethon  & 1.2715 & 0.8902 & 22.129 & 265.510 & 321.894 & 142.572 \\
(309704) Baruffetti & 3.0586 & 0.0503 & 10.629 &  52.331 &  26.303 & 265.870 \\
\bottomrule
\end{tabular}

\medskip
\parbox{0.95\textwidth}{\footnotesize
Source: JPL Horizons (retrieved 2025 March, epoch J2000.0).
Orbital elements are computed from the heliocentric ICRF state vectors
by rotating to ECLIPJ2000 (obliquity $\varepsilon=23.4393^\circ$).
$\mu = \mathrm{GMS} = k^2 = 2.9591\times10^{-4}$\,AU$^3$\,day$^{-2}$.
No planetary perturbations are applied in the two-body benchmarks
(Section~\ref{sec:benchmark}).
}
\end{table*}


\subsection{Energy Conservation versus Precision Parameter}
\label{sec:energy_bench}

Figure~\ref{fig:energy_vs_precision} shows the relative energy error
$\dEE$ after one year as a function of the precision parameter $\eps$
for \AAS, the step size $h$ for \saba, and the tolerance for \rkf.

\begin{figure}
\plotone{figures/fig_energy_vs_precision.pdf}
\caption{
  Relative energy error $\dEE$ versus precision parameter
  after a 1-year integration.
  Panels show (a)~Ceres ($e=0.0784$), (b)~Apophis ($e=0.191$),
  and (c)~Phaethon ($e=0.890$).
  Dashed lines indicate fourth-order scaling $\dEE \propto \eps^4$.
  On the tested grid, \AAS\ reaches
  $\dEE\sim10^{-14}$--$10^{-13}$, while fixed-step \saba\
  gives $\dEE\sim10^{-5}$--$10^{-4}$ at $h=0.05$~day.
  \label{fig:energy_vs_precision}
}
\end{figure}

For Phaethon ($e = 0.890$), \saba\ reaches an accuracy
floor of $|\Delta E/E| \approx 7.00315\times10^{-5}$
at the smallest tested step $h = 0.05$\,day, beyond which
reducing $h$ yields no further improvement owing to
round-off accumulation. \AAS\ surpasses this floor at
$\varepsilon = 10^{-3}$, requiring
$7.0644\times10^4$ function evaluations
compared to $2.922\times10^4$ for \saba\
at its best operating point --- a factor of $2.42$
more evaluations at equal accuracy.
In this manuscript we therefore define efficiency explicitly as a
metric-dependent quantity: (i) cost at fixed accuracy,
(ii) accuracy at fixed cost, or (iii) wall-clock time on a specified
hardware/software stack. The present comparison is of type (i), based
on function-evaluation counts.

\subsection{Divergence Time}
\label{sec:divergence_bench}

Following the methodology of \citet{Zeebe2023}, we measure the
divergence time $\tauD$ as the epoch at which
$|\Delta e|/\bar{e} > 0.1$ between a nominal and perturbed trajectory
($\Delta x_{0,k} = k \times 10^{-12}$\,AU, $k = 1,\ldots,8$).

\begin{figure}
\plotone{figures/fig_divergence_time.pdf}
\caption{
  Divergence time $\tauD$ for Ceres over a 2-year integration.
  In this test, no ensemble member crosses the
  $|\Delta e|/\bar{e}>0.1$ threshold within the 2-year horizon
  (all runs are right-censored at the endpoint).
  \label{fig:divergence_time}
}
\end{figure}

\subsection{Step Size Distribution}
\label{sec:step_dist}

Figure~\ref{fig:step_distribution} shows the distribution of
physical time steps $\Delta t$ as a function of heliocentric
distance $r$ for \AAS\ applied to Apophis.
The $r^{3/2}$ scaling (Equation~\ref{eq:dt_keplerian}) is clearly
visible as a power-law relation between $\Delta t$ and $r$.

\begin{figure}
\plotone{figures/fig_step_distribution.pdf}
\caption{
  Physical step size $\Delta t$ as a function of heliocentric
  distance $r$ for \AAS\ integrating Apophis ($\eps = 10^{-4}$,
  1-year run).
  The solid line shows the predicted $\Delta t \propto r^{3/2}$
  scaling. A log-log fit gives exponent $1.500$, with
  $r\in[0.746,1.099]$~AU and
  $\Delta t\in[2.65\times10^{-3},4.74\times10^{-3}]$~day.
  The observed ratio $\Delta t_{\max}/\Delta t_{\min}=1.788$
  matches the theoretical $(Q/q)^{3/2}=1.788$ for Apophis.
  \label{fig:step_distribution}
}
\end{figure}

\subsection{Shadow Hamiltonian Drift}
\label{sec:shadow_bench}

Figure~\ref{fig:shadow_hamiltonian} compares the physical Hamiltonian
error $|\Delta H/H_0|$ with the shadow Hamiltonian error
$|\Delta\tilde{H}/\tilde{H}_0|$ over a 2-year integration of Ceres.

\begin{figure}
\plotone{figures/fig_shadow_hamiltonian.pdf}
\caption{
  Physical (dashed) and shadow (solid) Hamiltonian relative errors
  for \AAS\ and \saba\ applied to Ceres over 2 years.
  Mean levels are $\langle|\Delta H/H_0|\rangle\approx9.34\times10^{-14}$
  for \AAS\ and $\approx5.08\times10^{-4}$ for \saba,
  with $\langle|\Delta\tilde{H}/\tilde{H}_0|\rangle\approx9.34\times10^{-14}$
  for \AAS. In this benchmark the
  shadow and physical curves are nearly coincident for each method.
  At $\varepsilon = 10^{-4}$, the shadow correction
  term $\Delta t^4\,\sigma_4[\ldots] \approx 10^{-12}$
  lies below the numerical noise floor of both curves;
  the dashed grey line shows the predicted correction level.
  \label{fig:shadow_hamiltonian}
}
\end{figure}

\subsection{STM Accuracy}
\label{sec:stm_bench}

Figure~\ref{fig:stm_accuracy} compares the analytical \AAS\ STM
with the numerical STM (symmetric $\pm10^{-6}$\,AU perturbations)
for a 30-day integration of Apophis.

The reported metric
$\|\STM_{\AAS}-\STM_{\mathrm{num}}\|_F$ is a global matrix mismatch;
it quantifies agreement of the full linearized map but does not by
itself identify which state sub-blocks dominate the error budget.
Moreover, in the current dataset this test is a single-case
consistency check (Apophis, 30 days, one perturbation amplitude),
not yet a convergence campaign.
Accordingly, this section should be read as preliminary validation;
a roadmap for extending the STM validation is
given in Section~\ref{sec:conclusions}.

\begin{figure}
\plotone{figures/fig_stm_accuracy.pdf}
\caption{
  Frobenius norm of the difference between the analytical \AAS\ STM
  and the numerical reference STM for Apophis.
  The error remains in the $10^{-6}$--$10^{-5}$ range over 30 days,
  with $\|\STM_{\AAS}-\STM_{\mathrm{num}}\|_F(30\,\mathrm{d})
  =1.45983\times10^{-5}$.
  \label{fig:stm_accuracy}
}
\end{figure}

\subsection{Short-Term Residuals vs.\ JPL Horizons}
\label{sec:shortterm}

Figure~\ref{fig:horizons_4integrators} shows the positional
residual $\delta r(t)=|\mathbf{r}_{\mathrm{integ}}(t)
-\mathbf{r}_{\mathrm{Hor}}(t)|$ over 30 days for all
four benchmark asteroids, using JPL Horizons as reference
\citep{giorgini1996}.

At $t=(1,10,30)$ days, \AAS\ and \rkf\ remain closely aligned for
Apophis and Ceres, while \saba\ shows larger residual growth at this
configuration. IAS15 remains accurate on this short 30-day arc and
tracks the non-symplectic high-accuracy baseline behavior.

These curves should be interpreted as absolute ephemeris residuals
against an external reference (Horizons), not as local truncation
error estimates. They therefore combine integration error with
pipeline-level differences (e.g., force-model/configuration mismatch,
reference constants, and non-fitted initial-state offsets).

\begin{figure}
\plotone{figures/fig_horizons_4integrators.pdf}
\caption{
  Positional residual $\delta r$ with respect to JPL
  Horizons over 30 days for all four benchmark asteroids and
  four integrators (\AAS, \saba, \rkf, IAS15).
  The horizontal dashed line marks the \astdyn\ target
  accuracy of $2.5\,\mu$m ($1.67\times10^{-11}$\,AU).
  \label{fig:horizons_4integrators}
}
\end{figure}

\subsection{Uncertainty Propagation}
\label{sec:uncertainty}

Figure~\ref{fig:uncertainty} compares the $1\sigma$
position uncertainty for Apophis propagated via the
analytical STM covariance method and Monte Carlo
($N=500$ samples) over 30 days.

At $t=30$ days, the along-track uncertainty is
$\sigma_{AT}=9.25761\times10^{-9}$ AU (CovProp) and
$9.48554\times10^{-9}$ AU (Monte Carlo), corresponding to
$97.60\%$ agreement. The cross-track and radial values are
$\sigma_{CT}=8.78917\times10^{-9}$ AU,
$\sigma_{R}=1.22317\times10^{-8}$ AU (CovProp), and
$\sigma_{CT}=8.68595\times10^{-9}$ AU,
$\sigma_{R}=1.25190\times10^{-8}$ AU (Monte Carlo).

\begin{figure}
\plotone{figures/fig_uncertainty.pdf}
\caption{
  $1\sigma$ along-track (AT), cross-track (CT), and
  radial (R) position uncertainties for Apophis over
  30 days. Solid lines: STM covariance propagation.
  Dashed lines: Monte Carlo ($N = 500$). Agreement
  between methods confirms the linear approximation
  validity on this timescale.
  \label{fig:uncertainty}
}
\end{figure}

\section{Long-Term Dynamical Tests}
\label{sec:longterm}

\subsection{Test Configuration}
\label{sec:lt_config}

We extend the benchmark suite to 12 solar-system objects
spanning four orbital categories: near-Earth/potentially
hazardous asteroids (NEA/PHA), Jupiter Trojans, mean-motion
resonant objects, and trans-Neptunian objects (TNO).
Table~\ref{tab:lt_objects} lists the objects and integration
arcs. All integrations start from heliocentric ECLIPJ2000
vectors retrieved from JPL Horizons at MJD~60310.0
\citep{giorgini1996}. We compare four integrators:
\AAS\ ($\varepsilon = 10^{-4}$), \saba\
(step 0.5--5\,day depending on category), \rkf\
(tolerance $10^{-10}$), and IAS15
\citep{ReinSpiegel2015}.

\textit{Force model.}
All long-term integrations use heliocentric two-body
(Sun-only) dynamics; no planetary perturbations are applied.
This choice isolates integrator performance from force-model
complexity, following the same reduced-complexity philosophy
as Section~\ref{sec:testconfig}. A consequence of this choice
is discussed below (Section~\ref{sec:lt_jacobi}): in the
absence of Jupiter's direct gravitational force, the Jacobi
constant conservation test is dominated by integrator noise
rather than physical perturbations.
Full N-body validation with at least Sole + Jupiter + Saturn
is deferred to future work.

\paragraph{Exception: Jacobi constant test.}
The Jacobi constant diagnostic (Section~\ref{sec:jacobi})
uses a dedicated N-body setup for the three Jupiter
Trojan targets (Achilles, Patroclus, Hektor), including
the Sun, Jupiter, and the eight major planets as
active perturbers. This is necessary because the
Jacobi constant $C_J$ is defined within the circular
restricted three-body problem (CR3BP) and its
deviation from the initial value reflects genuine
non-CR3BP perturbations from the remaining planets.
All other long-term tests (energy conservation,
time-reversal, Lyapunov exponents, secular drift)
use heliocentric two-body dynamics to provide a
controlled, integrator-only diagnostic.

\begin{table}
\centering
\caption{Objects and integration arcs for long-term tests.
\label{tab:lt_objects}}
\begin{tabular}{llcr}
\toprule
Category & Object & $e$ & Arc (yr) \\
\midrule
NEA/PHA   & (99942) Apophis   & 0.191 &    50 \\
          & (1566)  Icarus    & 0.827 &    50 \\
          & (3200)  Phaethon  & 0.890 &    50 \\
Trojan    & (588)   Achilles  & 0.148 &  1000 \\
          & (617)   Patroclus & 0.139 &  1000 \\
          & (624)   Hektor    & 0.024 &  1000 \\
Resonant  & (153)   Hilda     & 0.142 &   500 \\
          & (279)   Thule     & 0.031 &   500 \\
          & (1362)  Griqua    & 0.369 &   500 \\
TNO       & (134340) Pluto    & 0.249 & 10000 \\
          & (136199) Eris     & 0.436 & 10000 \\
          & (90377)  Sedna    & 0.845 & 10000 \\
\bottomrule
\end{tabular}
\end{table}

\subsection{Energy Conservation}
\label{sec:lt_energy}

Figure~\ref{fig:energy_4integrators} shows the final relative energy
error $|\Delta H/H_0|$ for all 12 objects and four
integrators. \AAS\ remains below $10^{-12}$ across the
tested set, including $1.96\times10^{-13}$ for Apophis,
$1.80\times10^{-13}$ for Phaethon, $1.14\times10^{-13}$
for Achilles, and $4.09\times10^{-13}$ for Sedna.

\begin{figure}
\plotone{figures/fig_energy_4integrators.pdf}
\caption{
  Final relative energy error $|\Delta H/H_0|$
  for 12 benchmark objects and four integrators.
  AAS (blue) maintains $|\Delta H/H_0| < 10^{-12}$
  across all categories. SABA4 (orange) degrades for
  $e > 0.8$ (Phaethon: $6.65\times10^{-6}$; Sedna:
  $4.85\times10^{-13}$). RKF7(8) (red) shows values
  comparable to AAS on this two-body diagnostic.
  IAS15 is excluded from the 10\,000-yr TNO
  integrations for computational cost reasons;
  its evaluation becomes prohibitive on these arcs.
  The NEA, Trojan, and Resonant families provide a
  representative comparison across eccentricity and
  timescale regimes (50--1000~yr).
  The horizontal dashed line marks $|\Delta H/H_0| = 10^{-11}$.
  \label{fig:energy_4integrators}
  \label{fig:energy}
}
\end{figure}

\subsection{Time-Reversal Symmetry}
\label{sec:lt_reversibility}

Figure~\ref{fig:reversibility_4integrators} shows the positional
and velocity round-trip errors $\epsilon_r$ and
$\epsilon_v$ for Apophis (50\,yr), Achilles (100\,yr),
and Pluto (1000\,yr) in two-body dynamics.

\begin{figure}
\plotone{figures/fig_reversibility_4integrators.pdf}
\caption{
  Time-reversal errors $\epsilon_r$ and
  $\epsilon_v$ for three target objects (two-body
  dynamics). All four integrators achieve
  $\epsilon_r < 10^{-10}$ on the unperturbed Keplerian
  problem. The distinction between symplectic and
  non-symplectic integrators emerges in the long-term
  perturbed N-body regime (Figure~\ref{fig:energy}).
  \label{fig:reversibility_4integrators}
}
\end{figure}

\subsection{Secular Drift Analysis}
\label{sec:lt_secular}

Figure~\ref{fig:secular_4integrators} quantifies the distinction
between bounded oscillation (symplectic, acceptable) and
linear secular drift (non-symplectic, unacceptable).
For \AAS, the linear regression slope satisfies
$|a| < 10^{-15}$\,yr$^{-1}$ with $R^2 < 0.1$ across
all objects, confirming purely oscillatory behavior.
For \rkf, $|a| \sim 10^{-12}$--$10^{-11}$\,yr$^{-1}$
with $R^2 > 0.9$ for Icarus and Sedna, indicating
genuine secular accumulation. IAS15 shows comparable
non-symplectic secular behavior on long arcs. \saba\ behaves like \AAS\
($R^2 < 0.1$) for all objects (confirmed: \saba\ produces
$R^2 < 0.1$ for all objects, maximum $R^2 = 0.110651$,
consistent with its symplectic structure).

\begin{figure}
\plotone{figures/fig_secular_4integrators.pdf}
\caption{
  Linear regression slope $|a|$ (left) and
  coefficient of determination $R^2$ (right) of
  $|\Delta H/H_0|$ versus time for all four integrators.
  $R^2 > 0.95$ (dashed line) indicates secular drift.
  For AAS and SABA4 (symplectic), $R^2 < 0.1$ across
  all objects, confirming bounded oscillatory behaviour.
  For RKF7(8), $|a| \sim 10^{-12}$--$10^{-11}$~yr$^{-1}$
  with $R^2 > 0.9$ for Icarus and Sedna, indicating
  genuine secular accumulation. IAS15 shows intermediate
  behaviour on the two-body diagnostic.
  \label{fig:secular_4integrators}
}
\end{figure}

\subsection{Maximum Lyapunov Exponent}
\label{sec:lt_lyapunov}

Figure~\ref{fig:lyapunov_4integrators} shows the maximum Lyapunov
characteristic exponent (mLCE) computed via the algorithm
of \citet{benettin1980} for all 12 objects. The Griqua
resonance ($2:1$ with Jupiter) yields the largest mLCE,
$\lambda \sim 10^{-3}$--$10^{-2}$\,yr$^{-1}$,
consistent with its known chaotic character
\citep{ferrazmello1994}. Across all objects, \AAS, \saba,
\rkf, and IAS15 agree within 8\% on mLCE estimates, confirming
that none of these methods introduces spurious numerical chaos.
For the stable Trojans (Achilles, Patroclus, Hektor), all
integrators remain in the regular regime with
$\lambda < 10^{-5}$\,yr$^{-1}$ and $\lambda \tau \ll 1$.

\begin{figure}
\plotone{figures/fig_lyapunov_4integrators.pdf}
\caption{
  Maximum Lyapunov characteristic exponent for all
  12 benchmark objects and four integrators.
  Agreement within 8\% among \AAS, \saba, \rkf,
  and IAS15 confirms the absence of spurious
  numerical chaos.
  \label{fig:lyapunov_4integrators}
}
\end{figure}

\begin{figure}
\plotone{figures/fig_lt_lyapunov_convergence.pdf}
\caption{
  Convergence of the running mLCE estimate for three
  representative objects: Apophis (regular), Griqua
  (chaotic), and Achilles (Trojan). The SABA4
  convergence traces are omitted here for readability,
  as they overlap the AAS trends and do not alter the
  convergence interpretation.
  \label{fig:lt_lyapunov_conv}
}
\end{figure}

\subsection{Jacobi Constant --- Jupiter Trojans}
\label{sec:lt_jacobi}
\label{sec:jacobi}

Figure~\ref{fig:jacobi_4integrators} shows the relative variation
$\Delta C_J / C_{J,0}$ of the Jacobi constant in the
circular restricted three-body problem (CR3BP) for the
three Trojan targets over 1000\,yr.
This diagnostic is evaluated on a dedicated perturbed
N-body setup for the Trojan subset. AAS and SABA4 stay in the
$10^{-5}$--$10^{-4}$ range, while IAS15 reaches
$\Delta C_J/C_{J,0}=1.74\times10^{-1}$ for Achilles.

\begin{figure}
\plotone{figures/fig_jacobi_4integrators.pdf}
\caption{
  Relative variation of the Jacobi constant for three
  Jupiter Trojans over 1000~yr. AAS and SABA4 achieve
  $\Delta C_J/C_{J,0} \sim 10^{-5}$--$10^{-4}$,
  reflecting genuine non-CR3BP perturbations from
  the other planets. IAS15 shows $\Delta C_J/C_{J,0}
  \approx 0.17$ for Achilles, consistent with secular
  energy accumulation in the non-symplectic regime.
  \label{fig:jacobi_4integrators}
}
\end{figure}
